\pdfoutput=1
\documentclass[11pt]{amsart}
\usepackage{amsmath,amssymb,amsthm,mathtools}
\usepackage[margin=1.2in]{geometry}
\usepackage{booktabs}
\usepackage{xcolor}
\usepackage[colorlinks=true,linkcolor=blue,citecolor=blue,urlcolor=blue]{hyperref}
\hypersetup{pdfauthor={Bernd Johannes Wuebben},
pdftitle={Relative metric estimates near reducible Yang--Mills connections}}
\setlength{\emergencystretch}{2em}
\newtheorem{theorem}{Theorem}[section]
\newtheorem{proposition}[theorem]{Proposition}
\newtheorem{lemma}[theorem]{Lemma}
\newtheorem{corollary}[theorem]{Corollary}
\theoremstyle{definition}
\newtheorem{definition}[theorem]{Definition}
\newtheorem{example}[theorem]{Example}
\theoremstyle{remark}
\newtheorem{remark}[theorem]{Remark}
\newcommand{\R}{\mathbb R}
\newcommand{\C}{\mathbb C}
\newcommand{\Hh}{\mathbb H}
\newcommand{\g}{\mathfrak g}
\newcommand{\kk}{\mathfrak k}
\newcommand{\cH}{\mathcal H}
\newcommand{\cG}{\mathcal G}
\newcommand{\ad}{\operatorname{ad}}
\newcommand{\im}{\operatorname{im}}
\newcommand{\Stab}{\operatorname{Stab}}
\newcommand{\tr}{\operatorname{tr}}
\newcommand{\id}{\operatorname{id}}
\newcommand{\Vol}{\operatorname{Vol}}
\newcommand{\Sp}{\operatorname{Sp}}
\newcommand{\diag}{\operatorname{diag}}
\DeclareMathOperator{\End}{End}
\DeclareMathOperator{\Hom}{Hom}
\DeclareMathOperator{\rank}{rank}
\DeclareMathOperator{\Rea}{Re}
\newcommand{\ip}[2]{\langle #1,#2\rangle}
\newcommand{\norm}[1]{\lVert #1\rVert}
\title[Relative metrics near reducible Yang--Mills connections]{Relative metric estimates near reducible \\Yang--Mills connections}
\author{Bernd Johannes Wuebben}
\date{October 28, 2026}
\subjclass[2020]{53D20 (primary); 58D27, 53C26, 58J50 (secondary)}
\keywords{Yang--Mills connections, singular symplectic reduction,
quotient metric, stabilizer, small eigenvalues, quaternionic Kähler quotient}

\begin{document}
\begin{abstract}
We study the natural $L^2$ metric near a central Yang--Mills connection
on a closed surface. In the classical quadratic symplectic model, we
prove a relative metric estimate with error quadratic in the harmonic
deformation parameter, uniformly across all stabilizer types. We also
compute the leading metric correction using the Green operators of the
reference connection. The proof first eliminates the gauge directions
with a fixed positive spectral gap and then minimizes over the remaining
finite-dimensional orbit; it requires no lower bound on a positive
orbit eigenvalue. The same argument compares every small eigenvalue of
the gauge Laplacian with the infinitesimal stabilizer-action operator,
including at changes of kernel dimension. As a curved finite-dimensional
comparison, quaternionic projective and hyperbolic reductions admit
exact orbit-metric formulas on their flat moment-map zero cone. A
four-cycle is realized in both settings with independently vanishing
weights. The resulting estimates distinguish uniform control of the
quotient metric from estimates that depend on the smallest orbit scale.
\end{abstract}
\maketitle

\section{Introduction}

Near a reducible connection, the orbit directions of the gauge group
contain vectors of arbitrarily small norm. This makes it difficult to
compare quotient metrics by estimates for an inverse gauge Laplacian:
the inverse becomes unbounded as the stabilizer grows. A finite-dimensional
model records the small directions, but an additive error estimate for
that model need not control its smallest positive eigenvalues.

For central Yang--Mills connections on a surface, the local symplectic
model is classical. Atiyah and Bott identify curvature as the moment
map for the gauge action \cite{AB83}. Huebschmann constructs a
stabilizer-equivariant Kuranishi chart that identifies the moduli germ
with a quadratic symplectic reduction \cite{Hue95,Hue96}. We use this
chart to compare the actual $L^2$ quotient metric with the metric of the
linear reduction. The estimates are relative and their constants do not
depend on the stabilizer type or on ratios between positive orbit
eigenvalues.

\subsection{The surface result}
Let $P\to\Sigma$ be a principal bundle with compact structure group $G$
over a closed oriented Riemannian surface. Fix an invariant positive
inner product on $\g$ and a central Yang--Mills connection $A_0$.
Thus $F_{A_0}=\zeta\,\mathrm{vol}_\Sigma$ for a parallel central
section $\zeta$. We consider connections with this fixed curvature,
modulo the gauge group $\cG$. All spaces and inner products
below are real. Put
\[
 d=d_{A_0},\qquad K=\Stab(A_0),\qquad
 \kk=\cH^0_d,\qquad V=\cH^1_d.
\]
Centrality gives $d^2=0$ on the adjoint bundle. The finite spaces
$\kk$ and $V$ have their $L^2$ inner products. Let $P_j$ be harmonic
projection in degree $j$, and let $Q_j=1-P_j$. Pairing with harmonic
zero-forms identifies $\cH^2_d$ with $\kk^*$. The quadratic moment map
and the infinitesimal action are
\begin{equation}\label{eq:intromodel}
 \Theta(h)=\tfrac12P_2[h,h],\qquad
 T_h:\kk\longrightarrow V,\quad T_h\xi=[h,\xi].
\end{equation}
An overall sign in the moment map depends on the convention for the
gauge action and does not change its zero set. Write $Z=\Theta^{-1}(0)$,
and let $g_{\mathrm{lin}}$ be the reduced metric on each orbit-type
stratum of $Z/K$.

The classical chart has the form $A(h)=A_0+a(h)$, with
$a(h)=h+c(h)$ and $c(h)\in\im d^*$. We recall its construction in
Section~\ref{sec:chart}. Its induced map on the quotient is denoted
$\overline a$.

\begin{theorem}\label{thm:metric}
There are $\delta>0$ and $C<\infty$ such that, for every
$h\in Z\cap B_\delta(V)$, the classical chart identifies the local
orbit-type strata and satisfies
\begin{equation}\label{eq:metricbound}
 (1-C|h|^2)g_{\mathrm{lin}}
 \ \leq\ \overline a^{\,*}g_{L^2}
 \ \leq\ (1+C|h|^2)g_{\mathrm{lin}}.
\end{equation}
The same constant applies to all strata. There is no restriction on
the ratios of the nonzero eigenvalues of $T_h^*T_h$.

The leading correction is explicit. Let
\[
 R_2=d^*G_2Q_2,\qquad
 \Delta_0=d^*d|_{\kk^\perp},\qquad
 B_hv=Q_0\ad_h^*v,
\]
where $G_2$ is the degree-two Green operator and
$\ad_h\eta=[h,\eta]$. For a tangent vector $v$ to the appropriate
smooth orbit-type piece of $Z$, chosen $L^2$-orthogonal to
$\im T_h$, one has
\begin{equation}\label{eq:metricexpansion}
 \begin{split}
 \overline a^{\,*}g_{L^2}([v],[v])
 ={}&|v|^2+\norm{R_2[h,v]}^2
       -\ip{B_hv}{\Delta_0^{-1}B_hv}+\mathcal R_h(v),\\
 &|\mathcal R_h(v)|\leq C|h|^3|v|^2.
 \end{split}
\end{equation}
This remainder is uniform across the orbit types as well.
\end{theorem}

Here a tangent vector means a tangent vector to a smooth piece, rather
than an arbitrary vector in $\ker d\Theta_h$ at a singular point.
The positive correction in \eqref{eq:metricexpansion} measures the
bending of the Coulomb slice. The negative correction is the part
removed by nonconstant gauge transformations. These terms use only
Green operators at $A_0$, where the relevant complementary gap is
fixed.

The same construction controls the small gauge spectrum. List
eigenvalues with multiplicities, including zero eigenvalues.

\begin{theorem}\label{thm:spectrum}
Let $m=\dim\kk$ and let
$0\leq\nu_1(h)\leq\cdots\leq\nu_m(h)$ be the eigenvalues of
$T_h^*T_h$. For $h\in Z\cap B_\delta(V)$, the self-adjoint gauge
Laplacian $d_{A(h)}^*d_{A(h)}$ has exactly $m$ eigenvalues below a fixed
positive threshold. They satisfy
\begin{equation}\label{eq:small}
 (1-C|h|^2)\nu_j(h)\leq\lambda_j(A(h))
       \leq(1+C|h|^2)\nu_j(h),\qquad 1\leq j\leq m.
\end{equation}
The next eigenvalue is bounded below uniformly, and
$\ker d_{A(h)}=\ker T_h\subset\kk$. The estimates include parameters
at which this kernel increases in dimension.
\end{theorem}

For unitary bundles, central Yang--Mills connections on a Riemann
surface are Hermitian--Einstein connections. The theorems therefore
compare actual gauge-theoretic metrics and operators. In higher complex
dimension, a quadratic tangent model need not describe the integrable
deformation locus; no such replacement is made here.

\subsection{Proof mechanism and finite-dimensional comparisons}
The gauge Lie algebra splits as $\kk\oplus\kk^\perp$. On
$\kk^\perp$, the covariant derivative at $A_0$ has a positive gap.
If $u$ is tangent to the Coulomb graph, then
\[
 \ip{u}{d_{A(h)}\eta}=\ip{u}{[a(h),\eta]},
 \qquad \eta\in\kk^\perp.
\]
The angle between the graph tangent and these gauge directions is
therefore $O(|h|)$. Orthogonal projection changes squared norms by
$O(|h|^2)$. Equivariance then identifies the remaining gauge action
exactly with $T_h$. Minimizing over its image preserves the relative
estimate, even if its dimension changes. A second expansion gives
\eqref{eq:metricexpansion}. The spectral argument uses the same
Coulomb cancellation in a Schur complement.

In finite dimensions, an equivariant Darboux chart at a Kähler fixed
point gives a related quadratic metric comparison simultaneously on
all local central moment levels. We record this consequence of the
classical construction in Section~\ref{sec:kahler}. It separates the
elementary passage from ambient to quotient estimates from the gauge
analysis needed in Theorem~\ref{thm:metric}.

There is also an explicit curved quaternionic comparison. For a compact
subgroup $K\subset\Sp(n)$ fixing the origin in the affine chart of
$\Hh P^n$ or $\Hh H^n$, the quaternionic moment-map zero set is the
same as the flat hyperkähler zero cone. If $r=|q|$ and
$\varepsilon=1$ or $-1$, respectively, its orbit Gram form satisfies
\begin{equation}\label{eq:introqk}
 G_\varepsilon(q)=(1+\varepsilon r^2)^{-1}G_0(q).
\end{equation}
Section~\ref{sec:quaternionic} gives the exact reduced metric and its
uniform bounds without a local-freeness hypothesis. Section~\ref{sec:cycle}
realizes the same weighted four-cycle in these curved quotients and in
the surface theorem. This identifies a common operator on an explicit
parameter family; it does not identify the two moduli spaces.

\subsection{Relation to existing normal forms}
The local symplectic reduction and its orbit-type decomposition are
due to Sjamaar--Lerman \cite{SL91}; holomorphic slice results are
developed by Sjamaar \cite{Sja95}. In the surface problem, the exact
quadratic moment map and local quotient chart are Huebschmann's
results \cite[Theorems 2.24 and 2.32]{Hue95}. Our proof retains their
coexact construction and estimates the full gauge quotient metric,
including its leading correction. We make no claim of a new
topological or symplectic local model.

Quaternionic Kähler reduction originates with Galicki--Lawson
\cite{GL88}. Swann describes its relation to hyperkähler reduction
\cite{Swa90} and the geometry of singular reductions \cite{Swa97}.
Mayrand proves a stratification and a local
complex-symplectic quotient model for integrable hyperkähler actions
\cite[Theorems 1.2--1.4]{May22}. Those structural results do not,
by themselves, compare an arbitrary curved quotient metric with a
flat one uniformly across all orbit types. Our quaternionic formulas
use the explicit space-form metric.

\section{Relative estimates and quotient norms}\label{sec:linear}

Let a compact group $K$ act on a Riemannian manifold and preserve a
subset $Z$ partitioned into smooth orbit-type pieces. At $z$ write
$T_z\xi=X_\xi(z)$ for the infinitesimal action and
$W_z=\im T_z$. On a smooth piece, the norm of a quotient tangent
class is
\begin{equation}\label{eq:quotientnorm}
 |[v]|_g^2=\inf_{w\in W_z}g_z(v-w,v-w).
\end{equation}
This formula also applies when the stabilizer has positive dimension.
It uses the image of the action, not an inverse on its Lie algebra.

\begin{lemma}\label{lem:relative}
Suppose two metrics on the same space satisfy
$(1-\eta)g_0\leq g_1\leq(1+\eta)g_0$, with $0\leq\eta<1$.
Then their orbit Gram forms and their reduced metrics satisfy the
same inequalities. The assertion also holds after pullback by an
equivariant diffeomorphism preserving $Z$.
\end{lemma}
\begin{proof}
Evaluate the metric inequalities on $T_z\xi$ for the orbit forms.
Apply them to every $v-w$ and take the infimum for the quotient forms.
Equivariance identifies the vertical spaces under a diffeomorphism.
\end{proof}

The following slightly sharper observation will be used for the metric
expansion. All norms in it are Euclidean operator norms.

\begin{lemma}\label{lem:residual}
Let $H=I+E$ be a positive form on a Euclidean space, with
$\norm E\leq\eta<1$, and let $W$ be any subspace. If $v\perp W$,
then
\[
 \inf_{w\in W}H(v-w,v-w)
 =H(v,v)-\ip{P_WE v}{(I_W+P_WE|_W)^{-1}P_WE v}.
\]
The absolute value of the subtracted term is at most
$\eta^2(1-\eta)^{-1}|v|^2$, independently of $W$.
\end{lemma}
\begin{proof}
Differentiation in $w$ gives
$(I_W+P_WE|_W)w=P_WE v$. The operator on $W$ is bounded below
by $1-\eta$. Substitution gives the formula and the bound.
\end{proof}

\subsection{Length metrics}
We use the intrinsic length metric associated with the metrics on the
orbit-type strata. Admissible paths are finite concatenations of paths
smooth in a stratum on the interior of each segment, continuous at the
endpoints, and of finite length. Endpoints may lie on smaller strata.
We take the infimum of lengths and then the metric completion when
needed. All local comparisons below use corresponding neighborhoods
under the stated chart. Thus a coordinate ball is not implicitly
identified with a ball for the reduced metric.

If relative quadratic-form bounds hold with a common constant on all
strata, their square roots compare lengths of every admissible path.
Taking infima gives the corresponding distance bounds, and the same
bounds extend to completions. This elementary observation is the
meaning of the assertions across singular strata in this paper.

An additive operator estimate cannot substitute for a relative one.
For example,
\[
 G(r)=\diag(r^2,r^8),\qquad
 \widetilde G(r)=G(r)+\diag(0,r^4)
\]
satisfy $\norm{\widetilde G-G}=r^4$, while their smallest positive
eigenvalues have unbounded ratio as $r\to0$. The estimates below
retain the energy in each orbit direction.

\section{The classical surface chart}\label{sec:chart}

We give the analytic details of the chart because the metric estimate
depends on its particular transverse choice. The symplectic and
moment-map assertions in this section are the classical surface local
model \cite[Section 2]{Hue95}.

\subsection{The coexact equation}
Fix an integer $s>2$ and use Sobolev spaces $H^s$ of adjoint-valued
forms. The bracket is continuous in the required multiplication
norms. Hodge decomposition gives
\[
 \Omega^1=\cH^1_d\oplus\im d\oplus\im d^*,\qquad
 \Omega^2=\cH^2_d\oplus\im d.
\]
In the completed spaces these images are closed. The map
\[
 d:\im d^*\cap H^s\Omega^1\longrightarrow
       \im d\cap H^{s-1}\Omega^2
\]
is an isomorphism, with inverse $R_2=d^*G_2Q_2$.

For $h\in V$ and $c\in\im d^*\cap H^s$, solve
\begin{equation}\label{eq:coexact}
 Q_2\left(d(h+c)+\tfrac12[h+c,h+c]\right)=0.
\end{equation}
Its derivative in $c$ at $(0,0)$ is the preceding isomorphism.
The implicit-function theorem gives a unique small solution $c(h)$,
smooth in $h$, and $K$-equivariant. Its value
and first derivative vanish at zero. In particular,
\begin{equation}\label{eq:cest}
 \norm{c(h)}_{H^s}\leq C|h|^2,\qquad
 \norm{Dc(h)v}_{H^s}\leq C|h||v|.
\end{equation}
Elliptic regularity applied to
$c=-\tfrac12R_2[h+c,h+c]$ gives smooth spatial coefficients.
The same estimates hold in every fixed $C^j$ norm after choosing
the initial Sobolev index sufficiently large. We use this finite
smooth family throughout; the constants may depend on that index.

Set $a(h)=h+c(h)$ and $U_h=Da(h)$. Harmonic and coexact forms are
orthogonal, so
\begin{equation}\label{eq:graphmetric}
 \norm{U_hv}^2=|v|^2+\norm{Dc(h)v}^2,
 \qquad d^*a(h)=d^*U_hv=0.
\end{equation}
The homogeneous quadratic term of $c$ and its derivative are
\begin{equation}\label{eq:cquadratic}
 c_2(h)=-\tfrac12R_2[h,h],\qquad
 J_hv=Dc_2(h)v=-R_2[h,v].
\end{equation}
Taylor's theorem in this finite parameter space gives
$Dc(h)=J_h+O(|h|^2)$ in each fixed spatial norm.

\subsection{The exact moment map}
The affine space of connections has symplectic form
$\Omega(u,v)=\int_\Sigma\langle u\wedge v\rangle$, where the
angle brackets pair the adjoint-bundle factors, with the conventional
sign chosen to match the curvature moment map. Harmonic forms pair
trivially with coexact forms for $\Omega$. Moreover $\im d^*$ is
isotropic. Indeed, on a surface it equals $*\im d$, star preserves
$\Omega$, and
\[
 \int_\Sigma\langle d\xi\wedge d\eta\rangle=0
\]
by integration by parts and $d^2=0$. It follows that
\begin{equation}\label{eq:symplecticchart}
 a^*\Omega=\Omega|_V.
\end{equation}
The moment map of the residual $K$-action on the graph is
$P_2(F_{A(h)}-F_{A_0})$, paired with $\kk$. Equivariance and
\eqref{eq:symplecticchart} show that its differential equals that
of the quadratic moment map $\Theta$ in \eqref{eq:intromodel}.
Both maps vanish at zero, and hence
\begin{equation}\label{eq:exactmoment}
 P_2(F_{A(h)}-F_{A_0})=\Theta(h).
\end{equation}
Together with \eqref{eq:coexact}, this proves that $A(h)$ has
curvature $F_{A_0}$ exactly when $h\in Z$.

The gauge slice theorem identifies a small Coulomb slice modulo $K$
with a neighborhood of $[A_0]$ in the gauge quotient. Applied to
\eqref{eq:exactmoment}, this is Huebschmann's local quotient theorem
\cite[Theorem 2.32]{Hue95}; its orbit-type interpretation is given
in \cite{Hue96}. We use that result for the local identification.
The subsequent norm comparison concerns the quotient by the full
gauge group, rather than only the metric induced on the slice.

The fixed central-curvature equation is preserved by the full gauge
group, including when $G$ is disconnected. To see this, let
$\pi_{\mathfrak z}$ be projection onto the adjoint center bundle.
Its flat connection $d_{\mathfrak z}$ is independent of the choice
of connection on $P$, and
\[
 \pi_{\mathfrak z}(F_{A_0+b}-F_{A_0})
        =d_{\mathfrak z}\pi_{\mathfrak z}b.
\]
Gauge transformations commute with $d_{\mathfrak z}$ and, being
orthogonal, with its adjoint. Thus the curvature of a gauge transform
of $A_0$ and $F_{A_0}$ are harmonic center-valued two-forms with
exact difference. They agree by Hodge theory. Gauge covariance
then proves preservation of the fixed-curvature equation.

Finally, differentiating $K$-equivariance gives an identity needed
in both proofs:
\begin{equation}\label{eq:equivariance}
 d_{A(h)}\xi=[a(h),\xi]=U_hT_h\xi,
 \qquad\xi\in\kk.
\end{equation}
One may reverse both infinitesimal-action signs without changing
any norm or operator in the argument.

\section{The actual quotient metric}\label{sec:metric}

Write $D_h=d_{A(h)}=d+\ad_{a(h)}$. The metric on the moduli stratum
is obtained by taking the infimum over all infinitesimal gauge
transformations. We first take that infimum over $\kk^\perp$.

\subsection{The complementary projection}
The Poincaré estimate at $A_0$ gives
\[
 \norm\eta_{H^1}\leq C\norm{d\eta}_{L^2},
 \qquad \eta\in H^1\Omega^0\cap\kk^\perp.
\]
Since $\norm{a(h)}_{C^0}\leq C|h|$, decreasing $\delta$ gives
\begin{equation}\label{eq:complementgap}
 \norm{D_h\eta}\geq(1-C|h|)\norm{d\eta}
                  \geq c\norm\eta_{H^1}.
\end{equation}
Thus $L_h=D_h|_{\kk^\perp}$ has closed range in $L^2\Omega^1$.
Let $\Pi_h$ be the orthogonal projection onto that range. For
$u=U_hv$, Coulombness in \eqref{eq:graphmetric} gives
\[
 \ip{u}{L_h\eta}=\ip{u}{[a(h),\eta]}.
\]
Using \eqref{eq:complementgap}, we obtain
\begin{equation}\label{eq:angle}
 |\ip{u}{L_h\eta}|
 \leq C|h|\norm u\norm{L_h\eta},\qquad
 \norm{\Pi_hu}\leq C|h|\norm u.
\end{equation}
Pythagoras yields the quadratic relative estimate
\begin{equation}\label{eq:complementnorm}
 (1-C|h|^2)\norm{U_hv}^2
 \leq\inf_{\eta\in H^1\cap\kk^\perp}
       \norm{U_hv-D_h\eta}^2
 \leq\norm{U_hv}^2.
\end{equation}
The use of squared projection norms is essential for the quadratic
error.

\subsection{Proof of the relative estimate}
Let $v$ represent a tangent vector to a smooth orbit-type piece of
$Z$. Every $H^1$ gauge generator splits uniquely as $\xi+\eta$,
with $\xi\in\kk$ and $\eta\perp\kk$. By
\eqref{eq:equivariance}, the full quotient norm is
\begin{equation}\label{eq:fullquotient}
 \overline a^{\,*}g_{L^2}([v],[v])
 =\inf_{\xi\in\kk,\,\eta\in H^1\cap\kk^\perp}
    \norm{U_h(v-T_h\xi)-D_h\eta}^2.
\end{equation}
Using smooth or higher-Sobolev gauge generators gives the same infimum
by density. Elliptic regularity identifies this norm with the usual
$L^2$ metric on the smooth moduli stratum.

Apply \eqref{eq:complementnorm} and \eqref{eq:graphmetric} to
$v-T_h\xi$, for every $\xi$, before taking the infimum. This gives
\[
 (1-C|h|^2)\inf_\xi|v-T_h\xi|^2
 \leq\overline a^{\,*}g_{L^2}([v],[v])
 \leq(1+C|h|^2)\inf_\xi|v-T_h\xi|^2.
\]
The outer expressions are the finite reduced norm. This proves
\eqref{eq:metricbound}. The proof uses neither an inverse of $T_h$
nor a bound on the dimension of its image away from the origin.

\subsection{The quadratic correction}
The closed positive form $\norm{L_h\eta}^2$ on
$H^1\Omega^0\cap\kk^\perp$ defines a self-adjoint operator
$C_h=L_h^*L_h$. It has domain $H^2\Omega^0\cap\kk^\perp$ and
a uniform positive lower bound by \eqref{eq:complementgap}.
The domain assertion follows from elliptic regularity; the fixed
harmonic projection has smooth finite-dimensional range.

Define a positive form $H_h$ on all of $V$ by
\begin{equation}\label{eq:Hform}
 H_h(v,v)=\norm{U_hv}^2
   -\ip{L_h^*U_hv}{C_h^{-1}L_h^*U_hv}.
\end{equation}
This is exactly the middle expression in
\eqref{eq:complementnorm}. All adjoints in this formula use the fixed
$L^2$ metrics. Since $d^*U_h=0$,
\begin{equation}\label{eq:sourceexpansion}
 L_h^*U_hv=Q_0\ad_{a(h)}^*U_hv
          =B_hv+O(|h|^2)|v|
\end{equation}
in $L^2$. Also
\[
 C_h-\Delta_0:
 H^2\Omega^0\cap\kk^\perp\longrightarrow
 L^2\Omega^0\cap\kk^\perp
\]
has norm $O(|h|)$, because its first-order and zeroth-order
coefficients are controlled by the fixed $C^1$ norm of $a(h)$.
The resolvent identity, or the Neumann series for
$I+(C_h-\Delta_0)\Delta_0^{-1}$, gives
\begin{equation}\label{eq:resolvent}
 \norm{C_h^{-1}-\Delta_0^{-1}}_{L^2\to H^2}\leq C|h|.
\end{equation}
Equations \eqref{eq:graphmetric} and \eqref{eq:cquadratic} give
\[
 \norm{U_hv}^2
 =|v|^2+\norm{R_2[h,v]}^2+O(|h|^3)|v|^2.
\]
Substituting this and \eqref{eq:sourceexpansion}--\eqref{eq:resolvent}
into \eqref{eq:Hform} proves
\begin{equation}\label{eq:Heffective}
 H_h(v,v)=|v|^2+\norm{R_2[h,v]}^2
     -\ip{B_hv}{\Delta_0^{-1}B_hv}+O(|h|^3)|v|^2.
\end{equation}
Every remainder here is a uniform operator-norm estimate on $V$.

For the residual quotient, choose $v\perp W_h=\im T_h$. By
\eqref{eq:fullquotient}, the desired norm is
$\inf_{w\in W_h}H_h(v-w,v-w)$. We have
$H_h=I+E_h$, with $\norm{E_h}\leq C|h|^2$.
Lemma~\ref{lem:residual} shows that this final infimum differs from
$H_h(v,v)$ by at most $C|h|^4|v|^2$, uniformly in $W_h$.
Combining with \eqref{eq:Heffective} proves
\eqref{eq:metricexpansion} and completes Theorem~\ref{thm:metric}.

\begin{corollary}\label{cor:length}
Let $d_{\mathrm{lin},R}$ and $d_{L^2,R}$ be the intrinsic length
metrics on $(Z\cap B_R)/K$ and its image, with the convention of
Section~\ref{sec:linear}. For sufficiently small $R$,
\[
 \sqrt{1-CR^2}\,d_{\mathrm{lin},R}
 \leq\overline a^{\,*}d_{L^2,R}
 \leq\sqrt{1+CR^2}\,d_{\mathrm{lin},R}.
\]
The comparison extends to their metric completions.
\end{corollary}
\begin{proof}
On every stratum segment in the truncated domain, Theorem~\ref{thm:metric}
compares speeds by the displayed square-root factors. Integrate,
sum over segments, and take the infimum over corresponding paths.
The two Lipschitz bounds extend the chart and its inverse to
completions.
\end{proof}

\begin{remark}
The expansion is pointwise on horizontal stratum tangent vectors.
Differentiating the horizontal projection as the stabilizer changes
can introduce inverse small eigenvalues. In particular,
\eqref{eq:metricexpansion} is not a curvature estimate for the
reduced metric.
\end{remark}

\section{All small gauge eigenvalues}\label{sec:spectrum}

We prove Theorem~\ref{thm:spectrum}. Regard
$\mathcal L_h=D_h^*D_h$ as a nonnegative self-adjoint operator on
$L^2\Omega^0$, with form domain $H^1\Omega^0$ and operator domain
$H^2\Omega^0$. It has compact resolvent. Use the fixed decomposition
$L^2\Omega^0=\kk\oplus\kk^\perp$ and write, in quadratic-form
blocks,
\begin{equation}\label{eq:block}
 \mathcal L_h=
 \begin{pmatrix}\Gamma_h&\mathcal B_h^*\\
                 \mathcal B_h&C_h\end{pmatrix}.
\end{equation}
The complementary operator is the one in the previous section;
$C_h\geq\gamma>0$ uniformly. Since $\kk$ consists of smooth
parallel sections, $\mathcal B_h=Q_0\mathcal L_h|_\kk$ is a
bounded map from the finite-dimensional space $\kk$ to $L^2$.

For $\xi\in\kk$, \eqref{eq:equivariance} and
\eqref{eq:graphmetric} show that
\begin{equation}\label{eq:Gamma}
 \norm{T_h\xi}^2
 \leq\ip{\xi}{\Gamma_h\xi}
 \leq(1+C|h|^2)\norm{T_h\xi}^2.
\end{equation}
There is also an exact cancellation. As $d\xi=0$ and $d^*a=0$,
\[
 d^*[a,\xi]=[d^*a,\xi]=0.
\]
Consequently
\begin{equation}\label{eq:spectralcoupling}
 \mathcal L_h\xi=\ad_a^*D_h\xi,
 \qquad
 \norm{\mathcal B_h\xi}^2
 \leq C|h|^2\ip{\xi}{\Gamma_h\xi}.
\end{equation}
The right side is the energy of the particular orbit direction,
which may be much smaller than $|h|^2|\xi|^2$.

For $0\leq\lambda<\gamma/2$, form the finite Schur operator
\[
 S_h(\lambda)=\Gamma_h-\lambda I
       -\mathcal B_h^*(C_h-\lambda)^{-1}\mathcal B_h.
\]
Equation~\eqref{eq:spectralcoupling} gives
\begin{equation}\label{eq:Schur}
 (1-C|h|^2)\Gamma_h-\lambda I
 \leq S_h(\lambda)\leq\Gamma_h-\lambda I.
\end{equation}
These finite inequalities control the eigenvalue count of the full
operator. Indeed, for $\xi\in\kk$ and
$\eta\in H^1\cap\kk^\perp$, completing the square gives
\begin{equation}\label{eq:inertia}
 \begin{split}
 \ip{\xi+\eta}{(\mathcal L_h-\lambda)(\xi+\eta)}
 ={}&\ip{\xi}{S_h(\lambda)\xi}\\
 &+\norm{(C_h-\lambda)^{1/2}
   \bigl(\eta+(C_h-\lambda)^{-1}\mathcal B_h\xi\bigr)}^2.
 \end{split}
\end{equation}
The left side is interpreted as a closed quadratic form.
The added vector belongs to $H^2$, so the change of coordinates is
invertible on the form domain. The number of eigenvalues of
$\mathcal L_h$ strictly below $\lambda$ is therefore the number
of negative eigenvalues of $S_h(\lambda)$.

Restriction to the trial space $\kk$ puts the first $m$ eigenvalues
below $\norm{\Gamma_h}=O(|h|^2)$. The max--min principle applied
to the codimension-$m$ subspace $\kk^\perp$ gives
$\lambda_{m+1}\geq\gamma$. If $\gamma_j(h)$ are the ordered
eigenvalues of $\Gamma_h$, including zeros, \eqref{eq:Schur} and
\eqref{eq:inertia} imply
\[
 (1-C|h|^2)\gamma_j(h)
 \leq\lambda_j(A(h))\leq\gamma_j(h),\qquad j\leq m.
\]
For the lower bound, evaluate the inertia comparison at
$\lambda=(1-C|h|^2)\gamma_j(h)$; at most $j-1$ eigenvalues are
strictly smaller. A zero right side gives the bound by positivity.
The upper bound is the ordinary min--max restriction. Repeated
eigenvalues cause no change. Finally \eqref{eq:Gamma} compares
$\gamma_j$ with $\nu_j$, proving \eqref{eq:small}.

To identify the kernel as a subspace, suppose
$D_h(\xi+\eta)=0$. Estimate \eqref{eq:complementnorm} applies to
$U_hT_h\xi=D_h\xi$, so its distance from $\im L_h$ is positive
unless $D_h\xi=0$. Hence $T_h\xi=0$ by injectivity of $U_h$,
and \eqref{eq:complementgap} gives $\eta=0$. The converse follows
from \eqref{eq:equivariance}. This proves the exact kernel assertion.
If the adjoint bundle is zero, the spectral assertion is read as the
absence of any spectrum; otherwise the preceding enumeration applies.

\section{A quantitative Kähler fixed-point comparison}\label{sec:kahler}

For completeness, the finite-dimensional estimate can be obtained
directly from the usual equivariant Darboux construction. Its proof is
short; its role is to isolate the part of the comparison that needs
no gauge-theoretic analysis.

\begin{proposition}\label{prop:kahler}
Let a compact group $K$ act by holomorphic isometries on a Kähler
manifold $(M,g,\omega)$ and fix $p$. Set $V=T_pM$, with its flat
metric $g_0$ and isotropy action. There is an equivariant symplectic
chart $\Phi:B_\delta(V)\to M$, fixing $p$, such that
\begin{equation}\label{eq:kahlerambient}
 (1-C|z|^2)g_0\leq\Phi^*g\leq(1+C|z|^2)g_0.
\end{equation}
In this chart the centered local moment map is exactly the quadratic
moment map of the isotropy representation. One chart and constant
therefore compare the orbit forms and reduced metrics on every
level $\lambda\in(\kk^*)^K$ in the domain, across all orbit types.
\end{proposition}

\begin{proof}
Let $E=\exp_p$ on an invariant ball. Normal coordinates and
$\nabla\omega=0$ give
\[
 E^*g=g_0+O(|z|^2),\qquad
 E^*\omega=\omega_0+\beta,\qquad
 \beta=O(|z|^2),\quad D\beta=O(|z|).
\]
The radial homotopy primitive
\[
 \alpha_z(v)=\int_0^1t\,\beta_{tz}(z,v)\,dt
\]
is invariant and satisfies $d\alpha=\beta$,
$\alpha=O(|z|^3)$ and $D\alpha=O(|z|^2)$.
The forms $\omega_t=\omega_0+t\beta$ are uniformly nondegenerate
on one small ball. Solve
$\iota_{X_t}\omega_t=-\alpha$. Then
$X_t=O(|z|^3)$ and $DX_t=O(|z|^2)$, uniformly in $t\in[0,1]$.
The flow $\psi_t$ is defined on a smaller invariant ball for all
these times. Its equation and variational equation give
\[
 \psi_1(z)=z+O(|z|^3),\qquad D\psi_1(z)=I+O(|z|^2).
\]
Thus $\Phi=E\circ\psi_1$ is equivariant and symplectic and
satisfies \eqref{eq:kahlerambient}.

The centered local moment map exists on this contractible
neighborhood and can be normalized equivariantly under the full
compact group. Equivariance and the symplectic identity show that
its pullback and the quadratic moment map have identical
differentials. Their values at zero agree, proving exact equality.
The levels invariant under all of $K$, rather than only its identity
component, are the levels stated in the proposition. Equivariance
preserves their stabilizers. Lemma~\ref{lem:relative} gives the
orbit and quotient estimates with the same constant.
\end{proof}

For a fixed radius $R$, the associated intrinsic distance estimates
have factors $\sqrt{1\pm CR^2}$. The proposition is a quantitative
consequence of equivariant Darboux and singular reduction
\cite{SL91,Sja95}; its proof uses the whole ambient metric estimate.
It does not attempt to invert a moment-map derivative on a singular
level. A triplet of moment maps on a quaternionic manifold requires
a separate treatment, as in the explicit class below.

\section{Quaternionic space-form quotients}\label{sec:quaternionic}

We now give a curved class in which the three moment equations and
the quotient metric can both be computed exactly. These formulas are
consequences of the classical Hopf construction and quaternionic
reduction \cite{GL88,Swa90}. They provide a useful comparison with
the surface estimate, without asserting a normal form for a general
quaternionic Kähler metric.

Let $K\subset\Sp(n)$ act on quaternionic columns by left matrix
multiplication. The maps are right-$\Hh$-linear. Extend the action
to $\Hh^{n+1}$ by fixing the first coordinate, and use the affine
chart $[1:q]$. Write
\[
 \varepsilon=1\quad\hbox{on }\Hh P^n,\qquad
 \varepsilon=-1\quad\hbox{on }\Hh H^n,\qquad
 s=1+\varepsilon|q|^2.
\]
The hyperbolic chart has $|q|<1$. Normalize both metrics to be the
Euclidean metric at $q=0$. Here $q^*v=\sum_j\overline q_jv_j$
is quaternion-valued.

\subsection{Metric and moment section}
The quaternionic Hermitian form
\[
 \mathcal B_\varepsilon(z,w)
   =\varepsilon\overline z_0w_0+\sum_{j=1}^n\overline z_jw_j
\]
has $\mathcal B_\varepsilon(x,x)=\varepsilon$ on the local Hopf
section $x(q)=(1,q)/\sqrt{s}$. For $b=q^*v$ and $a=\Rea b$,
direct differentiation gives
\[
 \mathcal B_\varepsilon(x,dx(v))=\frac{\operatorname{Im}b}{s},
 \qquad
 \Rea\mathcal B_\varepsilon(dx(v),dx(v))
      =\frac{|v|^2}{s}-\frac{\varepsilon a^2}{s^2}.
\]
The vertical vectors are $xu$, for $u\in\operatorname{Im}\Hh$,
with squared norm $\varepsilon|u|^2$. Subtracting their component
from $dx(v)$ gives the horizontal lift
$dx(v)-\varepsilon x\operatorname{Im}(b)/s$. Its squared norm is
\begin{equation}\label{eq:qkmetric}
 g_\varepsilon(v,v)
      =\frac{|v|^2}{s}-\frac{\varepsilon|q^*v|^2}{s^2}.
\end{equation}

The principal $\Sp(1)$ connection is
$\theta=\varepsilon\mathcal B_\varepsilon(x,dx)$.
For $\xi\in\kk$, its value on the lifted fundamental field is
\[
 \theta(\widetilde X_\xi)=\frac{\varepsilon q^*\xi q}{s}.
\]
The connection is $K$-invariant. Cartan's formula expresses the
covariant derivative of this section as minus contraction of the
curvature. Since that curvature is a nonzero constant multiple of
the quaternionic Kähler two-form triple, the quaternionic moment
section in this frame is a fixed nonzero scalar multiple of
$q^*\xi q/s$; see \cite[Chapter 3, \S3.2]{Swa90}.
The first homogeneous coordinate is fixed, so there is no constant
term in the numerator. In particular the curved zero set is exactly
the flat hyperkähler zero cone
\begin{equation}\label{eq:qkzero}
 Z_\Hh=\{q:q^*\xi q=0\text{ for every }\xi\in\kk\}
\end{equation}
intersected with the chart domain. This description uses a local
frame of the rank-three quaternionic bundle; it does not make the
moment section a globally defined triple of real functions.

\subsection{Exact comparison on every stratum}
Set $T_q\xi=\xi q$, and use one fixed invariant inner product on
$\kk$ to define orbit Gram operators. The smooth orbit-type pieces
of the flat reduction carry their usual quotient metrics
\cite[Theorem 1.1 and \S2.5]{May22}. The same pieces and quotient
norms apply to the curved zero set. Indeed, an element fixing
$[1:q]$ fixes the normalized lift $x(q)$: its possible right
$\Sp(1)$ multiplier is forced to be $1$ by the first coordinate.
The stabilizer therefore acts trivially on the quaternionic
rank-three bundle. Its fixed locus is the chart of a quaternionic
linear subspace, and reduction on the exact stabilizer locus by
$N_K(H)/H$ gives the smooth quotient pieces. This is the trivial
quaternionic-isotropy case of singular quaternionic reduction
\cite[\S\S3--4]{Swa97}; arbitrary singular quaternionic reductions
need not have only quaternionic Kähler strata.

\begin{proposition}\label{prop:qk}
On $Z_\Hh$, the flat and curved horizontal spaces agree, and
\begin{equation}\label{eq:qkgram}
 G_\varepsilon(q)=s^{-1}G_0(q).
\end{equation}
On each quotient stratum, the quaternion-valued one-form
$\beta_q(v)=q^*v$ descends and satisfies
\begin{equation}\label{eq:qkreduced}
 g_{\varepsilon,\mathrm{red}}
     =s^{-1}g_{0,\mathrm{red}}-\varepsilon s^{-2}|\beta|^2.
\end{equation}
Consequently, with $r=|q|$,
\begin{align}
 (1+r^2)^{-2}g_{0,\mathrm{red}}
   &\leq g_{+,\mathrm{red}}
     \leq(1+r^2)^{-1}g_{0,\mathrm{red}},\label{eq:qkplus}\\
 (1-r^2)^{-1}g_{0,\mathrm{red}}
   &\leq g_{-,\mathrm{red}}
     \leq(1-r^2)^{-2}g_{0,\mathrm{red}}.\label{eq:qkminus}
\end{align}
No local-freeness assumption is required for these estimates.
\end{proposition}

\begin{proof}
The zero equation says $q^*T_q\xi=0$. Polarization of
\eqref{eq:qkmetric} therefore gives, for every ambient vector $v$,
\[
 g_\varepsilon(v,T_q\xi)=s^{-1}g_0(v,T_q\xi).
\]
This proves both the Gram identity and equality of the orthogonal
complements of the orbit spaces, including where their ranks change.
The one-form $\beta$ is invariant, because $K\subset\Sp(n)$,
and annihilates the vertical vectors on $Z_\Hh$. It thus descends
on each stratum. Evaluate \eqref{eq:qkmetric} on the common
horizontal representative to obtain \eqref{eq:qkreduced}.
Finally $|q^*v|^2\leq r^2|v|^2$ gives
\eqref{eq:qkplus} and \eqref{eq:qkminus}.
\end{proof}

For the intrinsic distances on corresponding truncations $r\leq R$,
with the path convention of Section~\ref{sec:linear}, these bounds give
\[
 (1+R^2)^{-1}d_0\leq d_+\leq d_0,
 \qquad
 d_0\leq d_-\leq(1-R^2)^{-1}d_0\quad(R<1).
\]
They persist after completion. These are statements about the
specified coordinate domains, not about unrestricted global distances.

\subsection{Abelian charges and supports}
For a torus, let $Q_e$ be the integral charge of its $e$th quaternionic
coordinate and use the action
$(T_q\xi)_e=i(Q_e\cdot\xi)q_e$. Its flat Gram matrix is
\begin{equation}\label{eq:charge}
 G_0(q)=\sum_e|q_e|^2Q_eQ_e^{\mathsf T}.
\end{equation}
At a zero-moment point, every eigenvalue of the curved Gram operator
is exactly $s^{-1}$ times its flat value. This includes all zero
eigenvalues and uses the same fixed inner product on the Lie algebra.
A gauge-boson mass interpretation additionally requires a specified
positive gauge kinetic form; \eqref{eq:charge} alone does not supply it.

Zero-moment points with independently chosen supports can be obtained
by taking two copies of each charge. In a pair, set the coordinates
equal to $(b_e,jb_e)$, with $b_e\geq0$ real. Their moments cancel,
since
\[
 \overline b_e i b_e+\overline{jb_e}i(jb_e)
       =b_e^2\{i+(-j)ij\}=0,
\]
whereas their Gram contribution is $2b_e^2Q_eQ_e^{\mathsf T}$.
Thus supports of charge \emph{pairs} may be chosen independently.
Arbitrary individual-coordinate supports in an undoubled
representation need not satisfy the moment equations.

\section{A four-cycle in both geometric settings}\label{sec:cycle}

The following example makes the independence from weak orbit scales
explicit. Its matrix algebra is elementary. The geometric point is
that the same matrix occurs on actual solution sets in both settings.

\subsection{The common matrix}
Give the edges $12,23,34,41$ of a four-cycle the weights $a,s,a,t$,
where $a>0$ and $0\leq s,t\leq a/4$. Its Laplacian is
\begin{equation}\label{eq:cyclematrix}
 L(a,s,t)=
 \begin{pmatrix}
 a+t&-a&0&-t\\
 -a&a+s&-s&0\\
 0&-s&a+s&-a\\
 -t&0&-a&a+t
 \end{pmatrix}.
\end{equation}
The reflection $(14)(23)$ decomposes $\R^4$ into two orthogonal
planes. In the bases $(e_1+e_4,e_2+e_3)/\sqrt2$ and
$(e_1-e_4,e_2-e_3)/\sqrt2$, its blocks are
\[
 \begin{pmatrix}a&-a\\-a&a\end{pmatrix},\qquad
 \begin{pmatrix}a+2t&-a\\-a&a+2s\end{pmatrix}.
\]
Hence the ordered eigenvalues on $\mathbf1^\perp$ are
\begin{equation}\label{eq:cycleeigen}
 \nu_-=a+s+t-\sqrt{a^2+(s-t)^2},\qquad
 2a,\qquad
 \nu_+=a+s+t+\sqrt{a^2+(s-t)^2}.
\end{equation}
The inequalities $\sqrt{a^2+x^2}\leq a+x^2/(2a)$ and
$|s-t|\leq s+t$ give
\begin{equation}\label{eq:weakscale}
 \left(1-\frac{s+t}{2a}\right)(s+t)
       \leq\nu_-\leq s+t.
\end{equation}
Thus $\nu_-=0$ precisely when $s=t=0$, and
$\nu_-\sim s+t$ whenever $(s+t)/a\to0$, with no condition on
$s/t$. The product of the three values is
\[
 2a\nu_-\nu_+=4a^2(s+t)+8ast.
\]
These identities follow from the displayed blocks.

\subsection{Surface Hermitian--Einstein connections}
Let $\Sigma$ have genus $g\geq2$ and area one. Choose pairwise
nonisomorphic flat unitary line bundles $L_1,\ldots,L_4$ whose
tensor product is trivial, and fix a parallel unitary trivialization
of that product. Such a choice is obtained from three generic flat
lines by taking the fourth to be their inverse product. The direct
sum of the flat connections is then an $\mathrm{SU}(4)$ connection
on $E=\bigoplus_iL_i$, to which the preceding theorems apply. Its
infinitesimal determinant-one stabilizer is the space of diagonal
matrices $i\diag(\xi_1,\ldots,\xi_4)$ with
$\sum\xi_i=0$. With the invariant matrix inner product
$-\tr(XY)$ and area one, this is Euclidean $\mathbf1^\perp$.

For $i\ne j$, the degree-zero line bundle
$\Hom(L_j,L_i)$ is nontrivial. It has no holomorphic sections,
so Riemann--Roch gives
\[
 \dim_\C H^{0,1}(\Sigma,\Hom(L_j,L_i))=g-1.
\]
The real harmonic deformation space identifies with the underlying
real space of
\[
 H^{0,1}(\Sigma,\End_0 E),
\]
where $\End_0 E$
denotes the trace-free complex endomorphism bundle.
For each edge, select one harmonic direction in each of
the two ordered off-diagonal summands. Normalize its complex
coordinate so that the induced real harmonic-one-form norm is
the standard norm. Distinct selected summands are orthogonal.

The two summands on the edge $ij$ have opposite torus weights
$e_i-e_j$ and $e_j-e_i$. Assign them equal complex-coordinate
magnitudes $\sqrt{w_{ij}/2}$, where
$(w_{12},w_{23},w_{34},w_{41})=(a,s,a,t)$. Since the reference
stabilizer is diagonal, the moment map is precisely the torus
moment map on this orthogonal representation. The two
contributions cancel on every edge. We therefore obtain
\begin{equation}\label{eq:surfacecycle}
 h\in\Theta^{-1}(0),\qquad |h|^2=2a+s+t,\qquad
 T_h^*T_h=L(a,s,t)|_{\mathbf1^\perp}.
\end{equation}
For small weights the classical chart gives actual flat unitary
connections $A(h)$ with the specified determinant. They are
Hermitian--Einstein connections on their induced holomorphic
bundles. Their existence here follows from the exact moment
equation and the chart, not from the matrix calculation.

Theorem~\ref{thm:spectrum} compares the three small gauge
eigenvalues with \eqref{eq:cycleeigen} by factors
$1\pm C(2a+s+t)$. When $s+t>0$, the support graph is connected
and its determinant-one infinitesimal stabilizer vanishes. At
$s=t=0$ it has dimension one. The same estimate applies there,
including the newly zero eigenvalue. Theorem~\ref{thm:metric}
simultaneously compares the full quotient metrics, rather than
only this restricted family of orbit operators.

\subsection{Curved quaternionic quotients}
Let the effective torus $T^4/U(1)_{\mathrm{diag}}$ act on
$\Hh^8$. For each edge, take two quaternionic coordinates of
charge $Q_{ij}=e_i-e_j$, with values
\[
 \left(\sqrt{w_{ij}/2},\;j\sqrt{w_{ij}/2}\right).
\]
The paired-charge calculation gives
\[
 q\in Z_\Hh,\qquad |q|^2=2a+s+t,\qquad
 G_0(q)=L(a,s,t)|_{\mathbf1^\perp},
\]
using the Euclidean metric on the torus Lie algebra. The points
$[1:q]$ lie in the actual zero level of the quaternionic moment
section on $\Hh P^8$ and, when $2a+s+t<1$, on $\Hh H^8$.
Their curved orbit
eigenvalues are exactly
\begin{equation}\label{eq:curvedcycle}
 \frac{\nu_-}{1+\varepsilon(2a+s+t)},\qquad
 \frac{2a}{1+\varepsilon(2a+s+t)},\qquad
 \frac{\nu_+}{1+\varepsilon(2a+s+t)}.
\end{equation}
Equations \eqref{eq:surfacecycle} and \eqref{eq:curvedcycle}
compare actual geometric operators on a common parameter family.
The two quotient spaces have different dimensions and are not
identified by this construction.

For example, take $a=\rho^2$, $s=\rho^8$, and
$t=\exp(-1/\rho^2)$ as $\rho\downarrow0$. The smallest positive
model eigenvalue is asymptotic to $\rho^8$, whereas
$r^2=2a+s+t\sim2\rho^2$. An additive $O(r^4)$ error can
exceed that eigenvalue. Both geometric comparisons above instead
give relative error $O(\rho^2)$, regardless of the separation
between $s$ and $t$. Setting both weak weights to zero retains
the same estimates with the enlarged kernel.

\section{Further questions}

The surface argument uses two special features: the adjoint
connection at the reference point is flat, and the coexact graph
is exactly symplectic. In higher complex dimension, integrability
adds equations to the deformation problem. A useful extension
would require a chart of the actual integrable locus with a
uniform metric estimate; the orbit Gram matrix alone cannot
provide that chart.

The quaternionic space-form calculation makes a different
simplification. Its moment section and metric are explicit in
one common coordinate system. For a general quaternionic Kähler
metric, establishing a comparable chart across all stabilizer
types remains a separate problem. In particular, the formulas
above do not determine the metric of a physical compactification
whose metric is not already known.

Finally, the estimates control tangent norms and the resulting
local intrinsic distances. They do not control derivatives of
the quotient metric near a change of orbit type. Curvature
asymptotics would require additional estimates, since the
horizontal projection can vary rapidly even when the metric
comparison is uniform.

\begin{thebibliography}{Hue95}

\bibitem[AB83]{AB83}
M.~F. Atiyah and R.~Bott,
\emph{The Yang--Mills equations over Riemann surfaces},
Philos.\ Trans.\ Roy.\ Soc.\ London Ser.\ A \textbf{308} (1983), 523--615;
\href{https://doi.org/10.1098/rsta.1983.0017}{doi:10.1098/rsta.1983.0017}.

\bibitem[GL88]{GL88}
K.~Galicki and H.~B. Lawson, Jr.,
\emph{Quaternionic reduction and quaternionic orbifolds},
Math.\ Ann.\ \textbf{282} (1988), 1--21;
\href{https://doi.org/10.1007/BF01457009}{doi:10.1007/BF01457009}.

\bibitem[Hue95]{Hue95}
J.~Huebschmann,
\emph{The singularities of Yang--Mills connections for bundles on a
surface. I. The local model},
Math.\ Z.\ \textbf{220} (1995), 595--609;
\href{https://arxiv.org/abs/dg-ga/9411006}{arXiv:dg-ga/9411006}.

\bibitem[Hue96]{Hue96}
J.~Huebschmann,
\emph{The singularities of Yang--Mills connections for bundles on a
surface. II. The stratification},
Math.\ Z.\ \textbf{221} (1996), 83--92;
\href{https://arxiv.org/abs/dg-ga/9411007}{arXiv:dg-ga/9411007}.

\bibitem[May22]{May22}
M.~Mayrand,
\emph{Stratification of singular hyperkähler quotients},
Complex Manifolds \textbf{9} (2022), 261--284;
\href{https://doi.org/10.1515/coma-2021-0140}{doi:10.1515/coma-2021-0140}.

\bibitem[Sja95]{Sja95}
R.~Sjamaar,
\emph{Holomorphic slices, symplectic reduction and multiplicities of
representations},
Ann.\ of Math.\ (2) \textbf{141} (1995), 87--129;
\href{https://doi.org/10.2307/2118628}{doi:10.2307/2118628}.

\bibitem[SL91]{SL91}
R.~Sjamaar and E.~Lerman,
\emph{Stratified symplectic spaces and reduction},
Ann.\ of Math.\ (2) \textbf{134} (1991), 375--422;
\href{https://doi.org/10.2307/2944350}{doi:10.2307/2944350}.

\bibitem[Swa90]{Swa90}
A.~F. Swann,
\emph{HyperKähler and Quaternionic Kähler Geometry},
D.Phil.\ thesis, University of Oxford, 1990;
\url{https://users-math.au.dk/~swann/thesisafs.pdf}.

\bibitem[Swa97]{Swa97}
A.~Swann,
\emph{Singular moment maps and quaternionic geometry},
Banach Center Publ.\ \textbf{39} (1997), 143--153;
\url{https://matwbn.icm.edu.pl/ksiazki/bcp/bcp39/bcp39113.pdf}.

\end{thebibliography}

\bigskip
\noindent{\sc Bernd Johannes Wuebben, New York, NY,}
\texttt{wuebben@gmail.com}

\end{document}
